\par\Needspace{6\baselineskip}\medskip
\begin{center}\hypertarget{native-234-article169}{}\pdfbookmark[2]{169.}{article-169}169.\end{center}
Absoluto calculo art. praec. ante omnia arcum $C'C''$ determinare oportebit. Quod fiet commodissime, si antea perinde ut in art. 137 intersectio $D$ circulorum maximorum $A'C'B'$, $A''C''B''$ mutuaque inclinatio $\epsilon$ eruta fuerit: invenietur dein ex $\epsilon$, $C'D=z'+B'D$ atque $C''D=z''+B''D$, per formulas easdem quas in art. 144 tradidimus, non modo $C'C''=v''-v'$, sed etiam anguli $(u',u'')$, sub quibus circuli maximi $A'B'$, $A''B''$ circulum maximum $C'C''$ secant.

Postquam arcus $v''-v'$ inventus est, $v'-v$ et $r$ eruentur e combinatione aequationum
\[
r\sin(v'-v)=\frac{r''\sin(v''-v')}{P'}
\]
\[
r\sin(v'-v+v''-v')=\frac{1+P'}{P'}\cdot\frac{r'\sin(v''-v')}{1+\dfrac{Q'}{r'^3}}
\]
et perinde $r'''$ atque $v'''-v''$ e combinatione harum
\[
r'''\sin(v'''-v'')=\frac{r'\sin(v''-v')}{P''}
\]
\[
r'''\sin(v'''-v''+v''-v')=\frac{1+P''}{P''}\cdot\frac{r''\sin(v''-v')}{1+\dfrac{Q''}{r''^3}}\mathpunct{\cdot}
\]
\SourceLeaf{235}Omnes numeri hoc modo inventi exacti forent, si ab initio a valoribus veris ipsarum $P',P'',Q',Q''$ proficisci licuisset: tumque situm plani orbitae perinde ut in art. 149 vel ex $A'C',u'$ et $\gamma'$ vel ex $A''C'',u''$ et $\gamma''$ determinare conveniret, ipsasque orbitae dimensiones vel ex $r',r'',t',t''$ et $v''-v'$ vel, quod exactius est, ex $r,r''',t,t'''$ et $v'''-v$. Sed in calculo primo haec omnia praeteribimus, atque in id potissimum incumbemus, ut valores magis approximatos pro quantitatibus $P',P'',Q',Q''$ obtineamus. Hunc finem assequemur, si per methodum inde ab art. 88 expositam
\begin{center}
\begin{tabular}{@{}llll@{}}
ex & $r,\ r',\ v'-v,\ t'-t$ & eliciamus & $(\eta01)$\\
& $r',\ r'',\ v''-v',\ t''-t'$ & . . . . . . . . & $(\eta12)$\\
& $r'',\ r''',\ v'''-v'',\ t'''-t''$ & . . . . . . . . & $(\eta23)$.
\end{tabular}
\end{center}
Has quantitates, nec non valores ipsarum $r,r',r'',r''',\cos\tfrac12(v'-v)$ etc. in formulis III—VI substituemus, unde valores ipsarum $P',Q',P'',Q''$ resultabunt multo magis exacti quam ii, quibus hypothesis prima superstructa erat. Cum illis itaque hypothesis secunda formabitur, quae si prorsus eodem modo ut prima ad finem perducitur, valores ipsarum $P',Q',P'',Q''$ multo adhuc exactiores suppeditabit, atque sic ad hypothesin tertiam deducet. Hae operationes tam diu iterabuntur, donec valores ipsarum $P',Q',P'',Q''$ nulla amplius correctione opus habere videantur, quod recte iudicare exercitatio frequens mox docebit. Quoties motus heliocentricus parvus est, plerumque prima hypothesis illos valores iam satis exacte subministrat: si vero ille arcum maiorem complectitur, si insuper temporum intervalla ab aequalitate notabiliter recedunt, hypothesibus pluries repetitis opus erit; in tali vero casu hypotheses primae magnam calculi praecisionem haud postulant. In ultima denique hypothesi elementa ipsa ita ut modo indicavimus determinabuntur.
